\par\Needspace{9\baselineskip}\bigskip
\pdfbookmark[3]{Darstellung, [22]}{pallas-II-article22}
\begin{center}[22.]\end{center}
\begingroup
\setlength{\lineskip}{4pt}\setlength{\lineskiplimit}{4pt}
\noindent\hspace*{\parindent}Die momentanen Änderungen jedes Elements sind offenbar Produkte von \(\mathrm{d}t\) mit Funktionen der Zeit und der Elemente des \SourceLeaf{472}gestörten und des störenden Planeten. In der ersten Näherung werden alle diese Elemente als konstant angenommen. Die Zeit geht nur insofern ein, als sie in den Werten der mittleren Längen der beiden Planeten vorkommt, die sich ergeben zu
\[
\begin{aligned}
L&=N+nt\\
L'&=N'+n't,
\end{aligned}
\]
wobei wir alle auf den störenden Planeten bezogenen Größen durch Striche kennzeichnen. Auf diese Weise hat die momentane Änderung jedes Elements die Form \(\mathrm{d}t.(\text{Funktion von }L\text{ und }L')\). Es ist ferner ersichtlich, dass diese Funktion \textit{periodisch} von \(L\text{ und }L'\) abhängt, denn sie ändert ihren Wert [nicht], wenn man \(L\text{ und }L'\) um \(360^{0}\) oder um ein Vielfaches von \(360^{0}\) ändert. Funktionen dieser Art lassen sich bekanntlich durch eine Reihe von Gliedern der Form \(k\cos(iL+i'L')\) oder \(k\sin(iL+i'L')\) ausdrücken, wobei \(k\) ein konstanter Koeffizient ist, \(i\) alle positiven ganzen Zahlen (einschließlich null) und \(i'\) alle positiven oder negativen ganzen Zahlen durchläuft. Das Glied mit \(i=i'=0\) ergibt bei der Integration \(kt\), die säkulare Änderung des Elements. Die anderen Glieder ergeben periodische Teile der Form \(\dfrac{k}{in+i'n'}\sin(iL+i'L')\). Sollten sich Glieder mit \(in+i'n'=0\) finden, so wären sie auszunehmen und würden der säkularen Änderung \(kt\cos(iL+i'L')\) hinzufügen. Es ist jedoch ersichtlich, dass es solche Glieder nur geben kann, wenn die mittleren Bewegungen in einem rationalen Verhältnis stehen. Schließlich ist eine Konstante hinzuzufügen.
\par\endgroup
